% GENERATED FILE. Edit canonical lessons, then run npm run generate:books.
% canonical-source-hash: fnv1a64:dafba08cd8748f75
% canonical-lessons: PA-C226-tiar, PA-C226-burash, PA-C226-kambna, PA-C226-bherna, PA-C226-sair

\chapter{Ready for the Day}
\label{ch:pa-a2-ready-for-the-day}
\hlchaptermodality{226}

\begin{chapteropening}
\textbf{By the end of this chapter:} \emph{I can say how I get ready in the morning.}

Complete the last lesson of chapter 226: I can say how I get ready in the morning.
\end{chapteropening}

\section[\texorpdfstring{\pa{ਤਿਆਰ} (tiār)}{ਤਿਆਰ (tiār)}]{\pa{ਤਿਆਰ} (tiār) — ready}

\label{lesson:PA-C226-tiar}



\begin{quote}
Before the new one: say the Punjabi for to lie down, then the Punjabi for in a hurry.
\end{quote}

\subsection*{You'll want to know: \pa{ਤਿਆਰ}}
\textbf{\pa{ਤਿਆਰ}} — \emph{tiār} — \textquotedblleft{}ready\textquotedblright{}.

\begin{quote}
\pa{ਮੈਂ} \pa{ਸੱਤ} \pa{ਵਜੇ} \pa{ਤਿਆਰ} \pa{ਹੁੰਦੀ} \pa{ਹਾਂ।}
\end{quote}

\begin{quote}
\emph{maiṁ satt vaje tiār hundī hāṁ.}
\end{quote}

\begin{quote}
I am ready at seven.
\end{quote}

\begin{cousinweb}
\textbf{\pa{ਤਿਆਰ}} (\emph{tiār}) is borrowed: it is usually traced to Persian \emph{tayyār}, \textquotedblleft{}ready\textquotedblright{}.
\end{cousinweb}

\begin{grammarlens}[title={What you've built}]
One more word for saying what you do.
\end{grammarlens}

\subsection*{Your turn}
\begin{itemize}
\raggedright
  \item \emph{Say it:} \emph{tiār}
  \item \emph{Say it:} \emph{tiār}, once more
  \item \emph{Say it:} \emph{maiṁ satt vaje tiār hundī hāṁ.} — I am ready at seven.
  \item \emph{Recall:} say the Punjabi for to lie down, then the Punjabi for in a hurry, then say \emph{tiār} again
\end{itemize}

\begin{tcolorbox}[breakable,colback=teal!4,colframe=teal!35!black,title={Before you move on}]
What does \pa{ਤਿਆਰ} mean? (\textquotedblleft{}Ready\textquotedblright{}.) Say it once more.
\end{tcolorbox}
\section[\texorpdfstring{\pa{ਬੁਰਸ਼} (burash)}{ਬੁਰਸ਼ (burash)}]{\pa{ਬੁਰਸ਼} (burash) — a brush; a toothbrush}

\label{lesson:PA-C226-burash}



\begin{quote}
Before the new one: say the Punjabi for in a hurry, then the Punjabi for ready.
\end{quote}

\subsection*{You'll want to know: \pa{ਬੁਰਸ਼}}
\textbf{\pa{ਬੁਰਸ਼}} — \emph{burash} — \textquotedblleft{}a brush; a toothbrush\textquotedblright{}.

\begin{quote}
\pa{ਮੈਂ} \pa{ਸਵੇਰੇ} \pa{ਬੁਰਸ਼} \pa{ਨਾਲ} \pa{ਦੰਦ} \pa{ਸਾਫ਼} \pa{ਕਰਦਾ} \pa{ਹਾਂ।}
\end{quote}

\begin{quote}
\emph{maiṁ savere burash nāl dand sāf kardā hāṁ.}
\end{quote}

\begin{quote}
In the morning I clean my teeth with a brush.
\end{quote}

A masculine noun.

\begin{cousinweb}
\textbf{\pa{ਬੁਰਸ਼}} (\emph{burash}) is borrowed from English \emph{brush}.
\end{cousinweb}

\begin{grammarlens}[title={What you've built}]
One more word for saying what you do.
\end{grammarlens}

\subsection*{Your turn}
\begin{itemize}
\raggedright
  \item \emph{Say it:} \emph{burash}
  \item \emph{Say it:} \emph{burash}, once more
  \item \emph{Say it:} \emph{maiṁ savere burash nāl dand sāf kardā hāṁ.} — In the morning I clean my teeth with a brush.
  \item \emph{Recall:} say the Punjabi for in a hurry, then the Punjabi for ready, then say \emph{burash} again
  \item \emph{Write it:} copy \textbf{\pa{ਬੁਰਸ਼}} once, checking the whole word
\end{itemize}

\begin{tcolorbox}[breakable,colback=teal!4,colframe=teal!35!black,title={Before you move on}]
What does \pa{ਬੁਰਸ਼} mean? (\textquotedblleft{}A brush; a toothbrush\textquotedblright{}.) Say it once more.
\end{tcolorbox}
\section[\texorpdfstring{\pa{ਕੰਬਣਾ} (kambṇā)}{ਕੰਬਣਾ (kambṇā)}]{\pa{ਕੰਬਣਾ} (kambṇā) — to shiver}

\label{lesson:PA-C226-kambna}



\begin{quote}
Before the new one: say the Punjabi for ready, then the Punjabi for a toothbrush.
\end{quote}

\subsection*{You'll want to know: \pa{ਕੰਬਣਾ}}
\textbf{\pa{ਕੰਬਣਾ}} — \emph{kambṇā} — \textquotedblleft{}to shiver\textquotedblright{}.

\begin{quote}
\pa{ਸਰਦੀ} \pa{ਦੀ} \pa{ਸਵੇਰ} \pa{ਨੂੰ} \pa{ਮੈਂ} \pa{ਕੰਬਦਾ} \pa{ਹਾਂ।}
\end{quote}

\begin{quote}
\emph{sardī dī saver nū̃ maiṁ kambdā hāṁ.}
\end{quote}

\begin{quote}
On winter mornings I shiver.
\end{quote}

Past (I …): \pa{ਮੈਂ} \pa{ਬਹੁਤ} \pa{ਕੰਬਿਆ}, \emph{maiṁ bahut kambiā}.

\begin{cousinweb}
\textbf{\pa{ਕੰਬਣਾ}} (\emph{kambṇā}) is inherited from Sanskrit \emph{kampatē}, \textquotedblleft{}trembles, shivers\textquotedblright{}; CDIAL lists Punjabi \emph{kambṇā}, \textquotedblleft{}trembles, shivers\textquotedblright{} (CDIAL 2767).
\end{cousinweb}

\begin{grammarlens}[title={What you've built}]
One more word for saying what you do.
\end{grammarlens}

\subsection*{Your turn}
\begin{itemize}
\raggedright
  \item \emph{Say it:} \emph{kambṇā}
  \item \emph{Say it:} \emph{kambṇā}, once more
  \item \emph{Say it:} \emph{sardī dī saver nū̃ maiṁ kambdā hāṁ.} — On winter mornings I shiver.
  \item \emph{Recall:} say the Punjabi for ready, then the Punjabi for a toothbrush, then say \emph{kambṇā} again
  \item \emph{Write it:} copy \textbf{\pa{ਕੰਬਣਾ}} once, checking the whole word
\end{itemize}

\begin{tcolorbox}[breakable,colback=teal!4,colframe=teal!35!black,title={Before you move on}]
What does \pa{ਕੰਬਣਾ} mean? (\textquotedblleft{}To shiver\textquotedblright{}.) Say it once more.
\end{tcolorbox}
\section[\texorpdfstring{\pa{ਭੇੜਨਾ} (bheṛnā)}{ਭੇੜਨਾ (bheṛnā)}]{\pa{ਭੇੜਨਾ} (bheṛnā) — to shut a door}

\label{lesson:PA-C226-bherna}



\begin{quote}
Before the new one: say the Punjabi for a toothbrush, then the Punjabi for to shiver.
\end{quote}

\subsection*{You'll want to know: \pa{ਭੇੜਨਾ}}
\textbf{\pa{ਭੇੜਨਾ}} — \emph{bheṛnā} — \textquotedblleft{}to shut a door\textquotedblright{}.

\begin{quote}
\pa{ਮੈਂ} \pa{ਰਾਤ} \pa{ਨੂੰ} \pa{ਦਰਵਾਜ਼ਾ} \pa{ਭੇੜਦੀ} \pa{ਹਾਂ।}
\end{quote}

\begin{quote}
\emph{maiṁ rāt nū̃ darvāzā bheṛdī hāṁ.}
\end{quote}

\begin{quote}
At night I shut the door.
\end{quote}

Past (I …): \pa{ਮੈਂ} \pa{ਦਰਵਾਜ਼ਾ} \pa{ਭੇੜਿਆ}, \emph{maiṁ darvāzā bheṛiā}.

\begin{cousinweb}
\textbf{\pa{ਭੇੜਨਾ}} (\emph{bheṛnā}) is inherited from a reconstructed \emph{*bhiṭ}, \textquotedblleft{}meet, throng\textquotedblright{}; CDIAL lists Punjabi \emph{bheṛnā}, \textquotedblleft{}to close, compare\textquotedblright{} (CDIAL 9490).
\end{cousinweb}

\begin{grammarlens}[title={What you've built}]
One more word for saying what you do.
\end{grammarlens}

\subsection*{Your turn}
\begin{itemize}
\raggedright
  \item \emph{Say it:} \emph{bheṛnā}
  \item \emph{Say it:} \emph{bheṛnā}, once more
  \item \emph{Say it:} \emph{maiṁ rāt nū̃ darvāzā bheṛdī hāṁ.} — At night I shut the door.
  \item \emph{Recall:} say the Punjabi for a toothbrush, then the Punjabi for to shiver, then say \emph{bheṛnā} again
  \item \emph{Write it:} copy \textbf{\pa{ਭੇੜਨਾ}} once, checking the whole word
\end{itemize}

\begin{tcolorbox}[breakable,colback=teal!4,colframe=teal!35!black,title={Before you move on}]
What does \pa{ਭੇੜਨਾ} mean? (\textquotedblleft{}To shut a door\textquotedblright{}.) Say it once more.
\end{tcolorbox}
\section[\texorpdfstring{\pa{ਸੈਰ} (sair)}{ਸੈਰ (sair)}]{\pa{ਸੈਰ} (sair) — a walk; a stroll}

\label{lesson:PA-C226-sair}



\begin{quote}
Before the new one: say the Punjabi for to shiver, then the Punjabi for to shut a door.
\end{quote}

\subsection*{You'll want to know: \pa{ਸੈਰ}}
\textbf{\pa{ਸੈਰ}} — \emph{sair} — \textquotedblleft{}a walk; a stroll\textquotedblright{}.

\begin{quote}
\pa{ਮੈਂ} \pa{ਰੋਜ਼} \pa{ਸ਼ਾਮ} \pa{ਨੂੰ} \pa{ਪਿਤਾ} \pa{ਨਾਲ} \pa{ਸੈਰ} \pa{ਕਰਦਾ} \pa{ਹਾਂ।}
\end{quote}

\begin{quote}
\emph{maiṁ roz shām nū̃ pitā nāl sair kardā hāṁ.}
\end{quote}

\begin{quote}
Every evening I go for a walk with my father.
\end{quote}

A feminine noun.

\begin{cousinweb}
\textbf{\pa{ਸੈਰ}} (\emph{sair}) is borrowed: it is usually traced to Persian \emph{sair}, \textquotedblleft{}a walk, an outing\textquotedblright{}.
\end{cousinweb}

\begin{grammarlens}[title={What you've built}]
One more word for saying what you do.
\end{grammarlens}

\subsection*{Your turn}
\begin{itemize}
\raggedright
  \item \emph{Say it:} \emph{sair}
  \item \emph{Say it:} \emph{sair}, once more
  \item \emph{Say it:} \emph{maiṁ roz shām nū̃ pitā nāl sair kardā hāṁ.} — Every evening I go for a walk with my father.
  \item \emph{Recall:} say the Punjabi for to shiver, then the Punjabi for to shut a door, then say \emph{sair} again
  \item \emph{Write it:} copy \textbf{\pa{ਸੈਰ}} once, checking the whole word
\end{itemize}

\begin{tcolorbox}[breakable,colback=teal!4,colframe=teal!35!black,title={Before you move on}]
What does \pa{ਸੈਰ} mean? (\textquotedblleft{}A walk; a stroll\textquotedblright{}.) Say it once more.
\end{tcolorbox}
